\section{Introduction}

Ce mini-projet vise à concevoir et valider par simulation un système numérique sur puce articulé autour 
d'un processeur RISC-V et d'une mémoire modulaire, tout en mettant en application les principes de généricité 
et de réutilisation de blocs matériels (IP cores). La démarche consiste d'abord à transformer une mémoire 
RAM élémentaire en VHDL pour la rendre paramétrable et à instancier un bloc mémoire étendu (quadram) associé 
à un décodeur d'adresses. Dans un second temps, le modèle de processeur décrit en Verilog est adapté pour 
supporter l'extension d'instructions RV32IM et pour s'interfacer avec la mémoire externe. L'ensemble est 
enfin réuni au sein d'une entité englobante mixte VHDL/Verilog puis validé sous ModelSim à l'aide de programmes 
de test en assembleur, illustrant ainsi l'intégralité du flot de co-conception et d'intégration système.
\\

